% GENERATED FILE. Edit canonical lessons, then run npm run generate:books.
% canonical-source-hash: fnv1a64:c8052baae852dfe4
% canonical-lessons: SA-C455-sannidhyam, SA-C455-kamayate, SA-C455-mohah, SA-C455-asaktih, SA-C455-vatsalyam, SA-R455-review-448-455

\chapter{Fondness and Attachment}
\label{ch:sa-fondness-and-attachment}
\hlchaptermodality{455}

\begin{chapteropening}
\textbf{By the end of this chapter:} \emph{I can say what someone is fond of.}

Complete the last lesson of chapter 455: I can say what someone is fond of.
\end{chapteropening}

\section[\texorpdfstring{\sk{सान्निध्यम्} (sānnidhyam)}{सान्निध्यम् (sānnidhyam)}]{\sk{सान्निध्यम्} (sānnidhyam) — nearness; presence}

\label{lesson:SA-C455-sannidhyam}



\begin{quote}
Before the new one: say the Sanskrit for enmity, then the Sanskrit for covetous longing.
\end{quote}

\subsection*{You'll want to know: \sk{सान्निध्यम्}}
\textbf{\sk{सान्निध्यम्}} — \emph{sānnidhyam} — \textquotedblleft{}nearness; presence\textquotedblright{}.

\begin{quote}
\sk{मातुः} \sk{सान्निध्ये} \sk{शिशुः} \sk{शान्तः} \sk{भवति।}
\end{quote}

\begin{quote}
\emph{mātuḥ sānnidhye śiśuḥ śāntaḥ bhavati.}
\end{quote}

\begin{quote}
In its mother's presence the baby is calm.
\end{quote}

A neuter noun.

\begin{cousinweb}
\textbf{\sk{सान्निध्यम्}} (\emph{sānnidhyam}) is built on \emph{saṃnidhi}, \textquotedblleft{}nearness\textquotedblright{}, from \emph{sam}, \emph{ni} and the root √dhā, \textquotedblleft{}to put\textquotedblright{}, with the lengthening of its first vowel; MW gives \textquotedblleft{}the being near, nearness, vicinity, presence\textquotedblright{} (MW).
\end{cousinweb}

\begin{grammarlens}[title={What you've built}]
One more word for describing what you lived and felt.
\end{grammarlens}

\subsection*{Your turn}
\begin{itemize}
\raggedright
  \item \emph{Say it:} \emph{sānnidhyam}
  \item \emph{Say it:} \emph{sānnidhyam}, once more
  \item \emph{Say it:} \emph{mātuḥ sānnidhye śiśuḥ śāntaḥ bhavati.} — In its mother's presence the baby is calm.
  \item \emph{Recall:} say the Sanskrit for enmity, then the Sanskrit for covetous longing, then say \emph{sānnidhyam} again
\end{itemize}

\begin{tcolorbox}[breakable,colback=teal!4,colframe=teal!35!black,title={Before you move on}]
What does \sk{सान्निध्यम्} mean? (\textquotedblleft{}Nearness; presence\textquotedblright{}.) Say it once more.
\end{tcolorbox}
\section[\texorpdfstring{\sk{कामयते} (kāmayate)}{कामयते (kāmayate)}]{\sk{कामयते} (kāmayate) — desires}

\label{lesson:SA-C455-kamayate}



\begin{quote}
Before the new one: say the Sanskrit for covetous longing, then the Sanskrit for nearness.
\end{quote}

\subsection*{You'll want to know: \sk{कामयते}}
\textbf{\sk{कामयते}} — \emph{kāmayate} — \textquotedblleft{}desires\textquotedblright{}.

\begin{quote}
\sk{सर्वः} \sk{जनः} \sk{सुखं} \sk{कामयते।}
\end{quote}

\begin{quote}
\emph{sarvaḥ janaḥ sukhaṃ kāmayate.}
\end{quote}

\begin{quote}
Everyone desires happiness.
\end{quote}

Past (he …): \sk{सः} \sk{सुखं} \sk{कामितवान्}, \emph{saḥ sukhaṃ kāmitavān}.

\begin{cousinweb}
\textbf{\sk{कामयते}} (\emph{kāmayate}) comes from the root √kam, \textquotedblleft{}to love, to desire\textquotedblright{} (MW).
\end{cousinweb}

\begin{grammarlens}[title={What you've built}]
One more word for describing what you lived and felt.
\end{grammarlens}

\subsection*{Your turn}
\begin{itemize}
\raggedright
  \item \emph{Say it:} \emph{kāmayate}
  \item \emph{Say it:} \emph{kāmayate}, once more
  \item \emph{Say it:} \emph{sarvaḥ janaḥ sukhaṃ kāmayate.} — Everyone desires happiness.
  \item \emph{Recall:} say the Sanskrit for covetous longing, then the Sanskrit for nearness, then say \emph{kāmayate} again
\end{itemize}

\begin{tcolorbox}[breakable,colback=teal!4,colframe=teal!35!black,title={Before you move on}]
What does \sk{कामयते} mean? (\textquotedblleft{}Desires\textquotedblright{}.) Say it once more.
\end{tcolorbox}
\section[\texorpdfstring{\sk{मोहः} (mohaḥ)}{मोहः (mohaḥ)}]{\sk{मोहः} (mohaḥ) — infatuation; delusion}

\label{lesson:SA-C455-mohah}



\begin{quote}
Before the new one: say the Sanskrit for nearness, then the Sanskrit for desires.
\end{quote}

\subsection*{You'll want to know: \sk{मोहः}}
\textbf{\sk{मोहः}} — \emph{mohaḥ} — \textquotedblleft{}infatuation; delusion\textquotedblright{}.

\begin{quote}
\sk{धनस्य} \sk{मोहः} \sk{जनं} \sk{मूर्खं} \sk{करोति।}
\end{quote}

\begin{quote}
\emph{dhanasya mohaḥ janaṃ mūrkhaṃ karoti.}
\end{quote}

\begin{quote}
Infatuation with money makes a person foolish.
\end{quote}

A masculine noun.

\begin{cousinweb}
\textbf{\sk{मोहः}} (\emph{mohaḥ}) comes from the root √muh, \textquotedblleft{}to be bewildered\textquotedblright{}; MW gives \textquotedblleft{}bewilderment, infatuation, delusion\textquotedblright{} (MW).
\end{cousinweb}

\begin{grammarlens}[title={What you've built}]
One more word for describing what you lived and felt.
\end{grammarlens}

\subsection*{Your turn}
\begin{itemize}
\raggedright
  \item \emph{Say it:} \emph{mohaḥ}
  \item \emph{Say it:} \emph{mohaḥ}, once more
  \item \emph{Say it:} \emph{dhanasya mohaḥ janaṃ mūrkhaṃ karoti.} — Infatuation with money makes a person foolish.
  \item \emph{Recall:} say the Sanskrit for nearness, then the Sanskrit for desires, then say \emph{mohaḥ} again
\end{itemize}

\begin{tcolorbox}[breakable,colback=teal!4,colframe=teal!35!black,title={Before you move on}]
What does \sk{मोहः} mean? (\textquotedblleft{}Infatuation; delusion\textquotedblright{}.) Say it once more.
\end{tcolorbox}
\section[\texorpdfstring{\sk{आसक्तिः} (āsaktiḥ)}{आसक्तिः (āsaktiḥ)}]{\sk{आसक्तिः} (āsaktiḥ) — attachment}

\label{lesson:SA-C455-asaktih}



\begin{quote}
Before the new one: say the Sanskrit for desires, then the Sanskrit for infatuation.
\end{quote}

\subsection*{You'll want to know: \sk{आसक्तिः}}
\textbf{\sk{आसक्तिः}} — \emph{āsaktiḥ} — \textquotedblleft{}attachment\textquotedblright{}.

\begin{quote}
\sk{वृद्धस्य} \sk{गृहे} \sk{आसक्तिः} \sk{अस्ति।}
\end{quote}

\begin{quote}
\emph{vṛddhasya gṛhe āsaktiḥ asti.}
\end{quote}

\begin{quote}
The old man is attached to his house.
\end{quote}

A feminine noun.

\begin{cousinweb}
\textbf{\sk{आसक्तिः}} (\emph{āsaktiḥ}) is built from \emph{ā} and the root √sañj, \textquotedblleft{}to cling, to stick\textquotedblright{}; MW gives \textquotedblleft{}devotedness, attachment\textquotedblright{} (MW).
\end{cousinweb}

\begin{grammarlens}[title={What you've built}]
One more word for describing what you lived and felt.
\end{grammarlens}

\subsection*{Your turn}
\begin{itemize}
\raggedright
  \item \emph{Say it:} \emph{āsaktiḥ}
  \item \emph{Say it:} \emph{āsaktiḥ}, once more
  \item \emph{Say it:} \emph{vṛddhasya gṛhe āsaktiḥ asti.} — The old man is attached to his house.
  \item \emph{Recall:} say the Sanskrit for desires, then the Sanskrit for infatuation, then say \emph{āsaktiḥ} again
\end{itemize}

\begin{tcolorbox}[breakable,colback=teal!4,colframe=teal!35!black,title={Before you move on}]
What does \sk{आसक्तिः} mean? (\textquotedblleft{}Attachment\textquotedblright{}.) Say it once more.
\end{tcolorbox}
\section[\texorpdfstring{\sk{वात्सल्यम्} (vātsalyam)}{वात्सल्यम् (vātsalyam)}]{\sk{वात्सल्यम्} (vātsalyam) — tenderness for a child}

\label{lesson:SA-C455-vatsalyam}



\begin{quote}
Before the new one: say the Sanskrit for infatuation, then the Sanskrit for attachment.
\end{quote}

\subsection*{You'll want to know: \sk{वात्सल्यम्}}
\textbf{\sk{वात्सल्यम्}} — \emph{vātsalyam} — \textquotedblleft{}tenderness for a child\textquotedblright{}.

\begin{quote}
\sk{मातुः} \sk{वात्सल्यं} \sk{शिशुं} \sk{रक्षति।}
\end{quote}

\begin{quote}
\emph{mātuḥ vātsalyaṃ śiśuṃ rakṣati.}
\end{quote}

\begin{quote}
A mother's tenderness protects her baby.
\end{quote}

A neuter noun.

\begin{cousinweb}
\textbf{\sk{वात्सल्यम्}} (\emph{vātsalyam}) is built on \emph{vatsala}, \textquotedblleft{}fond, tender\textquotedblright{}, from \emph{vatsa}, \textquotedblleft{}a calf\textquotedblright{} (\textbf{\sk{वत्सः}}); MW gives \textquotedblleft{}affection or tenderness, esp. towards offspring\textquotedblright{} (MW).
\end{cousinweb}

\begin{grammarlens}[title={What you've built}]
One more word for describing what you lived and felt.
\end{grammarlens}

\subsection*{Your turn}
\begin{itemize}
\raggedright
  \item \emph{Say it:} \emph{vātsalyam}
  \item \emph{Say it:} \emph{vātsalyam}, once more
  \item \emph{Say it:} \emph{mātuḥ vātsalyaṃ śiśuṃ rakṣati.} — A mother's tenderness protects her baby.
  \item \emph{Recall:} say the Sanskrit for infatuation, then the Sanskrit for attachment, then say \emph{vātsalyam} again
\end{itemize}

\begin{tcolorbox}[breakable,colback=teal!4,colframe=teal!35!black,title={Before you move on}]
What does \sk{वात्सल्यम्} mean? (\textquotedblleft{}Tenderness for a child\textquotedblright{}.) Say it once more.
\end{tcolorbox}
\section[(dialogue)]{Looking back: from Naming a Feeling to Fondness and Attachment}

\label{lesson:SA-R455-review-448-455}



\begin{quote}
Say the Sanskrit for a feeling and regret.
\end{quote}

\subsection*{Your turn}
Say these aloud:

\begin{itemize}
\raggedright
  \item a feeling, regret, thrills with joy, merriment, distress
  \item a flurry of feeling, is amazed, a marvel, pleases, a thrill
  \item dejection, grieves, weeping, tears, a sigh
  \item is pained, takes heart, coaxes, troubled, anguish
  \item anxiety, despairing, loses heart, agitated, wretched
  \item pining for someone absent, bereft, lives away from home, a dear friend, sends word
  \item envies, detests, hatred, enmity, covetous longing
  \item nearness, desires, infatuation, attachment, tenderness for a child
\end{itemize}

\begin{tcolorbox}[breakable,colback=teal!4,colframe=teal!35!black,title={Before you move on}]
A feeling? (\textbf{\sk{भावः}}, \emph{bhāvaḥ}.) Regret? (\textbf{\sk{अनुतापः}}, \emph{anutāpaḥ}.) Thrills with joy? (\textbf{\sk{हृष्यति}}, \emph{hṛṣyati}.) Merriment? (\textbf{\sk{उल्लासः}}, \emph{ullāsaḥ}.) Distress? (\textbf{\sk{खेदः}}, \emph{khedaḥ}.) A flurry of feeling? (\textbf{\sk{आवेगः}}, \emph{āvegaḥ}.) Is amazed? (\textbf{\sk{विस्मयते}}, \emph{vismayate}.) A marvel? (\textbf{\sk{आश्चर्यम्}}, \emph{āścaryam}.) Pleases? (\textbf{\sk{तोषयति}}, \emph{toṣayati}.) A thrill? (\textbf{\sk{रोमाञ्चः}}, \emph{romāñcaḥ}.) Dejection? (\textbf{\sk{विषादः}}, \emph{viṣādaḥ}.) Grieves? (\textbf{\sk{शोचति}}, \emph{śocati}.) Weeping? (\textbf{\sk{रोदनम्}}, \emph{rodanam}.) Tears? (\textbf{\sk{बाष्पः}}, \emph{bāṣpaḥ}.) A sigh? (\textbf{\sk{निःश्वासः}}, \emph{niḥśvāsaḥ}.) Is pained? (\textbf{\sk{व्यथते}}, \emph{vyathate}.) Takes heart? (\textbf{\sk{आश्वसिति}}, \emph{āśvasiti}.) Coaxes? (\textbf{\sk{अनुनयति}}, \emph{anunayati}.) Troubled? (\textbf{\sk{व्याकुलः}}, \emph{vyākulaḥ}.) Anguish? (\textbf{\sk{सन्तापः}}, \emph{santāpaḥ}.) Anxiety? (\textbf{\sk{उद्वेगः}}, \emph{udvegaḥ}.) Despairing? (\textbf{\sk{हताशः}}, \emph{hatāśaḥ}.) Loses heart? (\textbf{\sk{विषीदति}}, \emph{viṣīdati}.) Agitated? (\textbf{\sk{व्यग्रः}}, \emph{vyagraḥ}.) Wretched? (\textbf{\sk{दीनः}}, \emph{dīnaḥ}.) Pining for someone absent? (\textbf{\sk{उत्कण्ठा}}, \emph{utkaṇṭhā}.) Bereft? (\textbf{\sk{विधुरः}}, \emph{vidhuraḥ}.) Lives away from home? (\textbf{\sk{प्रवसति}}, \emph{pravasati}.) A dear friend? (\textbf{\sk{सुहृद्}}, \emph{suhṛd}.) Sends word? (\textbf{\sk{सन्दिशति}}, \emph{sandiśati}.) Envies? (\textbf{\sk{ईर्ष्यति}}, \emph{īrṣyati}.) Detests? (\textbf{\sk{जुगुप्सते}}, \emph{jugupsate}.) Hatred? (\textbf{\sk{द्वेषः}}, \emph{dveṣaḥ}.) Enmity? (\textbf{\sk{वैरम्}}, \emph{vairam}.) Covetous longing? (\textbf{\sk{स्पृहा}}, \emph{spṛhā}.) Nearness? (\textbf{\sk{सान्निध्यम्}}, \emph{sānnidhyam}.) Desires? (\textbf{\sk{कामयते}}, \emph{kāmayate}.) Infatuation? (\textbf{\sk{मोहः}}, \emph{mohaḥ}.) Attachment? (\textbf{\sk{आसक्तिः}}, \emph{āsaktiḥ}.) Tenderness for a child? (\textbf{\sk{वात्सल्यम्}}, \emph{vātsalyam}.)
\end{tcolorbox}
